\documentclass[a4paper,11pt]{article}
\usepackage[T1]{fontenc}
\usepackage[utf8]{inputenc}
\usepackage{t1enc}
\usepackage[dvips]{color}
\usepackage{amsmath,amssymb,makeidx,stmaryrd}
\usepackage[normalem]{ulem}
\usepackage{fancybox}
\definecolor{gris}{gray}{0.9}
\usepackage{pst-plot,pst-tree,pstricks,pst-node,pst-text}
\usepackage{pst-eucl}
\usepackage{pstricks-add}
\usepackage[left=3.5cm, right=3.5cm, top=3cm, bottom=3cm]{geometry}
\newcommand{\vect}[1]{\overrightarrow{\,\mathstrut#1\,}}
\def\bas{\fcolorbox{gris}{gris}{\parbox{0.986\textwidth}{
	\color{red}\footnotesize 1995\hfill Page : \thepage/1}}}
\def\haut{\fcolorbox{gris}{gris}{\parbox{0.986\textwidth}{
	\color{red}
	\bf Concours Général de Mathématiques \hfill 1995
	}}}
\makeatletter
\def\ps@cg{%
\let\@oddhead\haut%
\let\@evenhead\haut%
\let\@oddfoot\bas
\let\@evenfoot\bas}
\makeatother
\newcommand{\R}{\mathbb{R}}
\newcommand{\N}{\mathbb{N}}
\newcommand{\D}{\mathbb{D}}
\newcommand{\Z}{\mathbb{Z}}
\newcommand{\Q}{\mathbb{Q}}
\newcommand{\C}{\mathbb{C}}
\pagestyle{cg}
\renewcommand{\theenumi}{\textbf{\arabic{enumi}}}
\renewcommand{\labelenumi}{\textbf{\theenumi.}}
\renewcommand{\theenumii}{\textbf{\alph{enumii}}}
\renewcommand{\labelenumii}{\textbf{\theenumii.}}
\parindent0pt
\everymath{\displaystyle}
\usepackage[french]{babel}

\DecimalMathComma
\usepackage[np]{numprint}
\begin{document}

\begin{center}\textbf{{\color{blue} Exercice I}}
\end{center}

Dans un plan $P$, on se donne un triangle ABC. 

À toute droite $D$, non parallèle à l'un de ses côtés, on associe le point G$_{\text{D}}$, isobarycentre des trois points communs à D et aux droites (AB), (BC) et (CA). 

L'objet de cet exercice est de déterminer l'ensemble $\mathcal{F}$ des points G$_{\text{D}}$ lorsque $D$ varie.

\medskip

\begin{enumerate}
\item Démontrer, que lorsque $D$ se déplace en restant parallèle à une droite $\delta$, le point G$_{\text{D}}$ décrit une droite $\Delta_{\delta}$.
\item On suppose, dans cette question seulement, ABC équilatéral. 

Montrer que, lorsque $\delta$ varie, les droites $\Delta_{\delta}$ sont toutes tangentes à un même cercle et déterminer l'ensemble $\mathcal{F}$ dans ce cas.
\item On revient au cas général. 

Montrer que l'on peut trouver un triangle équilatéral A$'$B$'$C$'$ de l'espace dont le projeté orthogonal sur le plan $P$ est le triangle ABC et en déduire l'ensemble $\mathcal{F}$.
\end{enumerate}

\begin{center}\textbf{{\color{blue} Exercice II}}
\end{center}

\medskip

Etudier la convergence de la suite $\left(u_n\right)_{n\in \N}$ définie par 

\[\left\{\begin{array}{l c l}
u_0		&\geqslant 	&0\: \text{et pour tout } n \in  \N :\\
u_{n+1}	& =			&\sqrt{u_n} + \dfrac{1}{n + 1}\\
\end{array}\right.\]

\begin{center}\textbf{{\color{blue} Exercice III}}
\end{center}

\medskip

Dans le plan, on considère $\Gamma_1, \: \Gamma_2$ et $\Gamma_3$ trois cercles de rayon $R$ passant par le point O,et on note $D$ l'ensemble des points du plan intérieurs à au moins deux de ces cercles.

Comment doit-on placer $\Gamma_1, \: \Gamma_2$ et $\Gamma_3$ pour que l'aire de $D$ soit minimale ? 

Justifier votre réponse.

\begin{center}\textbf{{\color{blue} Exercice IV}}
\end{center}

\medskip

Soient A$_1$, A$_2$, A$_3$, B$_1$, B$_2$, B$_3$ six points du plan tels que l'on ait pour tous les entiers $i$ et $j$ de $\{1,\: 2,\: 3\}$ :

\[\text{A}_i\text{B}_j = i + j.\]

Que peut-on dire de ces six points ?

\end{document}